% Source: Wald GR, physical PDF p. 27 (printed p. 14).
% Coverage: Figure 2.2, tangent plane at a point on a sphere.
% Figures/tables/footnotes: Figure only; no tables or footnotes.
% Status: source-reviewed and compile-clean.
% Uncertainties: none.
\begingroup
\renewcommand{\thefigure}{2.2}
\begin{figure}[htbp]
  \centering
\begin{tikzpicture}[x=0.78cm,y=0.78cm,>=Stealth,font=\small]
  \coordinate (P) at (0.52,0.73);
  % Put p on the near edge of the projected plane.  The plane lies on the
  % opposite side of that tangent line from the sphere, so it cannot be read
  % as a secant plane passing through the sphere.
  \draw (-0.89,-0.67) -- (2.57,2.76) -- (2.02,3.31) -- (-1.44,-0.12) -- cycle;
  \draw (1.42,-0.18) circle (1.28);
  \fill (P) circle (2pt);
  \node[below right] at (P) {\(p\)};
  \draw[->,thick] (P) -- ++(1.08,1.02);
\end{tikzpicture}
  \caption{The tangent plane at point \(p\) of a sphere in \(\mathbb{R}^3\).}
  \label{fig:2.2}
\end{figure}
\endgroup
